\scanpage{56}
% @ lines 1–3, top-edge return to earlier sentence / text
de $\unsure{\mathfrak g|g}=\operatorname{Lie}G$ $+$ \inserted{sont de la forme $\check{\mathbb V}(A^*)/G_m$ -- i.e. sont des ss-grpes ``\uline{jacobiens}''.} En effet, considérons $\mathfrak h=\operatorname{Lie}\unsure{\mathscr J|\mathscr T|\mathcal T}$, qui est un tel ss-fibré, ($\mathscr J$ lisse de dim rel. $1$)
% @ lines 4–6 / text
considérons le \deleted{$\widetilde{\mathscr J}$ qu'elle \unsure{déf.}} sous-grpe jacobien \unsure{qui il|qu'il} définit, qui est donc aussi le centralisateur de $\mathfrak h$ dans $G$.
% @ lines 6–7, left edge / margin note
% slanted
\margin{à vérifier ! \uline{il faut} \uuline{car. \unsure{rés.|p.} $\ne2$}}
% @ lines 7–10 / text
\deleted{Ceci signifie} Comme $\mathscr J$ est \uline{commutatif} il centralise $\mathfrak h:\operatorname{Lie}\mathscr J$, donc $\mathscr J\subset\widetilde{\mathscr J}$, d'où égalité puisque $\widetilde J$ est connexe et $J,\widetilde J$ sont lisses de $=$ dim rel. $1$\ldots
% archive sheet number: 55
